\chapter{Results}
The final result of the project is a functional Android application that transforms the educational activity of the \textit{Jardín de los Sentidos} into a cooperative, geolocated serious game. The application allows players to create a multiplayer session, join it using a code and take part in a shared experience in which the group explores the real garden and completes missions linked to its sensory areas.

The application includes the main systems required to support the intended experience. It integrates geolocation so that the player’s real position can be used inside the game, a map scene that acts as the central navigation space, a cooperative multiplayer structure that keeps the session synchronised between devices, and a set of missions connected to the senses represented in the garden. The game also includes feedback after each mission and a final screen showing the group name, the participants and the score achieved during the session.

From a gameplay perspective, the result is a complete playable flow. Players can start a session, access the map, move through the garden, activate missions, solve them cooperatively and obtain a final result. The experience is structured around exploration, observation, communication and group decision-making, maintaining the relationship with the physical environment throughout the game.

From an educational perspective, the application provides the lecturers with a digital tool that can be used to support the teaching session related to the \textit{Jardín de los Sentidos}. The final score and the group result can be used as part of the activity carried out in class, allowing the teachers to assess participation and performance during the session. In this sense, the project has produced a usable prototype that can be incorporated into the educational context for which it was designed.

\section{Project applications}
\textit{Micorrizas} has been developed as a tool that can be applied within the subject related to the natural environment in Early Childhood Education. Its main application is to complement the existing visit to the \textit{Jardín de los Sentidos} by providing a structured, interactive and cooperative activity that students can complete in the real space. If any problems arise, such as poor lighting because the activity is taking place late in the day or adverse weather conditions.

The game can be used by lecturers to guide students through the garden without requiring constant direct instruction. The application gives players objectives, activates quests according to their position and provides feedback on their answers. This allows the teaching activity to become more autonomous while preserving the importance of direct contact with the natural environment.

Another possible application is its use as an assessable class activity. Since the game records the group’s progression and provides a final score, the lecturers can use the result as part of the evaluation of the session. This does not replace the educational value of observation and discussion, but it provides an additional element that can help structure and assess the activity.

The project also creates a foundation for future educational uses. New plant species, sounds, questions, clues or missions could be added to adapt the experience to future editions of the subject or to different learning objectives. Therefore, the application is not limited to the current application, but can be expanded as a reusable educational resource.

